\chapter{Salsify: Low-Latency Network Video by Integrating the Codec and the Transport}
\paperinfo{S. Fouladi, J. Emmons, E. Orbay, C. Wu, R. S. Wahby, K. Winstein (Stanford University; Saratoga High School)}{USENIX NSDI 2018}

\begin{keybox}[The paper in three sentences]
Video-call systems (Skype, FaceTime, WebRTC) have a transport that estimates the network's \emph{average} rate and a separate video encoder that
aims its \emph{average} bitrate at that estimate, adjusting about once a second. So when the network suddenly slows down, oversized frames keep
going out and pile up in queues. Salsify merges the two into one per-frame loop. Before every frame, the transport says how many bytes the network
can take right now; a \key{purely functional} VP8 codec, whose state can be saved and reused, encodes the frame at two qualities from the same state;
Salsify sends whichever fits, or skips the frame entirely. The same saved-state trick lets it recover from packet loss by encoding from a state the
receiver has confirmed it holds.
\end{keybox}

% ------------------------------------------------------------------
\section{Background: why a video encoder has ``state''}

Salsify only makes sense once you see how video compression works, so start there. A video encoder does not compress each frame on its own.
Most frames are encoded as \textbf{differences from earlier frames}: ``this block is the same as that block in the previous picture, shifted 3
pixels left, plus this small correction''. So to decode frame 10, the decoder needs the pictures it decoded before. It keeps them as
\textbf{reference images}, together with probability tables used for the final compression step. Together these are the \key{codec state}, about
4\,MiB for 720p video.

The encoder keeps its own copy of the decoder's state, because it must know exactly what the decoder has in order to describe the next frame as a
difference. Two consequences follow:
\begin{itemize}
  \item \textbf{Encoder and decoder must stay in lock-step.} If the decoder misses a frame, its state no longer matches what the encoder assumed, and
        every later frame decodes as garbage until something resynchronises them.
  \item \textbf{A conventional encoder cannot take a frame back.} As soon as it encodes frame 10, it updates its internal state as if the decoder will
        receive frame 10. If the application then decides not to send that frame, the encoder's state is wrong. So the application can only decide
        \emph{before} encoding (``don't feed it this picture''), never after.
\end{itemize}

A second fact: \textbf{an encoder cannot hit an exact size}. How many bytes a frame compresses to depends on its content (a still face compresses
well, a turning head does not), and the encoder only learns the size after it has done the work. It can be told to aim for ``1\,Mb/s on average'',
but not ``make this frame 40\,KB''.

% ------------------------------------------------------------------
\section{The problem with today's design}

Today's real-time video systems are built from two separate boxes:
\begin{itemize}
  \item a \textbf{transport} that sends packets, reads ACKs and congestion signals, and estimates the network's \emph{average} rate;
  \item a \textbf{codec} with its own control loop, which chooses a frame rate and quality so that its \emph{average} bitrate matches that estimate,
        updated on roughly a one-second timescale (WebRTC, for example).
\end{itemize}

\textbf{What happens when the network suddenly slows down} (illustrative numbers): a cellular link drops from 10\,Mb/s to 2\,Mb/s. The encoder is
still producing 30 frames per second sized for 10\,Mb/s, about 40\,KB each. Until the next rate update (up to a second later), around 30 such frames,
1.2\,MB in total, are pushed into a link that now drains 0.25\,MB per second. They wait in the bottleneck queue for seconds, or overflow it and are
lost. The transport sends each oversized frame anyway and only \emph{afterwards} tells the codec to back off. The video freezes or lags badly, and
the delay was caused by the application's own traffic.

That is the regime Salsify targets: \textbf{Internet access links, especially cellular}, where
\begin{itemize}
  \item loss and queueing delay are caused mainly by how much \emph{you} send (you overflow the bottleneck buffer);
  \item the capacity can fall quickly;
  \item losses come in bursts, so forward error correction (sending redundant data) does not help much.
\end{itemize}
The classic research on video over networks assumed the opposite: a fixed rate and random, independent losses that the application must survive.
Salsify's goal is different: \textbf{never send more than the network can take right now}, so the loss and delay never happen. It still sends ordinary
UDP packets; it does not reach into the radio layer.

None of Salsify's pieces is new on its own: the VP8 format (2008), the functional VP8 codec (from the authors' ExCamera system), loss recovery in the
style of Mosh, and rate estimation from Sprout. The contribution is how they are wired together.

% ------------------------------------------------------------------
\section{The architecture: one loop, run once per frame}

\begin{center}
\begin{tabularx}{\linewidth}{@{}lXX@{}}
\toprule
 & \textbf{Conventional (WebRTC, Skype, \dots)} & \textbf{Salsify}\\
\midrule
Control loops & two: transport estimates rate, codec matches it & \textbf{one}, merged\\
What the codec aims at & an average bitrate over time & the size of \textbf{this} frame\\
How often & about once a second & \textbf{every frame}\\
Hitting the size & predict encoder settings beforehand & \textbf{encode twice, pick afterwards}\\
Frame rate & fixed; can only stop feeding pictures in & \textbf{variable}: can drop a frame after encoding it\\
Oversized frame & sent anyway, react later & never sent\\
Loss recovery & retransmit lost packets & encode from a state the receiver has confirmed\\
\bottomrule
\end{tabularx}
\end{center}

\textbf{The loop, every frame:}
\begin{enumerate}
  \item The \textbf{transport} works out how many bytes the network can accept right now without building a queue: \texttt{max\_frame\_size}.
  \item The \textbf{codec} encodes the new camera picture \textbf{twice}, from the \textbf{same starting state}: once at one quality step \emph{above}
        the last sent frame's quality, and once at one step \emph{below}.
  \item Salsify looks at the two actual sizes. It sends the \textbf{better} version if it fits, else the \textbf{worse} one if it fits, else
        \textbf{neither} (the frame is skipped).
  \item The state produced by whichever frame was sent becomes the starting state for the next frame. If the frame was skipped, the next frame starts from
        the same state as this one did.
\end{enumerate}
Because the two candidates bracket the last quality, quality climbs one step when there is room and falls one step when there is not, so it tracks capacity
frame by frame. Salsify-1c encodes the two versions one after the other on one core; Salsify-2c encodes them in parallel on two cores.

Back to the sudden slowdown: within one frame time the transport's budget shrinks. The next frame is sent at the lower quality if that fits,
or skipped if not. The queue never builds up.

% ------------------------------------------------------------------
\section{The purely functional codec}

The step ``encode twice and keep one'' is impossible for a normal encoder, which updates its hidden state on every encode. Salsify's VP8 codec
(about 11,000 lines of C++) makes the state an explicit value that is passed in and returned, with no side effects:
\begin{lstlisting}
decode(state, frame)          -> (new_state, image)
encode(state, image, quality) -> frame
\end{lstlisting}
The decoder is an automaton: given a starting state and a compressed frame, it outputs a picture and moves to a new state. Because nothing is hidden,
the application can:
\begin{itemize}
  \item \textbf{try several encodings from the same state} and throw away the ones it does not send. Nothing is disturbed, because the state was never
        modified in place;
  \item \textbf{skip a frame after encoding it}, just by not sending it and reusing the old state for the next frame. The frame rate therefore adapts to the
        network automatically;
  \item \textbf{choose which state to encode from}, which is the basis of loss recovery (next).
\end{itemize}
This also gets around the ``encoders cannot hit a size'' problem. Salsify does not need to predict the size, because it makes two attempts and
measures them. The decision is delayed until the exact sizes are known, which is the latest possible moment.

\subsection{Recovering from loss by choosing the starting state}
\begin{enumerate}
  \item \textbf{Normally} the sender encodes each frame from the state that the \emph{previous} frame leads to, without waiting for it to be
        acknowledged. It assumes frames in flight will arrive.
  \item The receiver acknowledges with a \textbf{hash of its current decoder state}, so the sender always knows the last state the receiver
        definitely has (\texttt{known\_receiver\_codec\_state}).
  \item \textbf{If a packet is lost}, the receiver cannot decode that frame and stays in the older state, so every later frame, built on the lost one, is
        useless to it.
  \item When ACKs reveal the loss, the sender enters \textbf{loss-recovery mode} for 5 seconds and encodes new frames from \texttt{known\_receiver\_codec\_state},
        a state the receiver has confirmed it holds. Those frames decode correctly, so video resumes with the very next frame, and nothing old is
        retransmitted.
\end{enumerate}

\textbf{Example.} The sender sends frames F1--F5, each built on the previous one. F3 is lost. The receiver decodes F1 and F2, so it is in state $S_2$,
and must discard F4 and F5 because they build on F3. ACKs carrying the hash of $S_2$ tell the sender that the receiver is stuck there. The sender encodes F6
from $S_2$ rather than from $S_5$. The receiver decodes F6 immediately and is back in sync. WebRTC would instead retransmit the lost packets of F3 and
wait for them, causing a delay spike after every outage before fresh video flows again.

To make this work, both sides \textbf{keep old states} until they are no longer needed. When the receiver gets a frame built on state $S$, it can discard all
states older than $S$. When the sender gets an ACK for $S$, it can do the same.

\textbf{Reordering.} If a fragment of a \emph{newer} frame arrives before the current frame is complete, the receiver gives up on the incomplete frame
(its state becomes invalid) and moves on. The sender sees this as a loss and recovers as above. Reordering \emph{within} one frame's fragments is
harmless, because the receiver waits for all of them before decoding.

% ------------------------------------------------------------------
\section{The video-aware transport}

\subsection{Packets and acknowledgements}
The transport (about 2,000 lines of C++) splits each frame into MTU-sized UDP fragments. Every fragment's header carries the frame number, the
fragment index, and \textbf{hashes of the frame's source state and target state}. A frame is therefore a self-describing instruction: ``from state X,
this produces state Y''. The receiver can tell immediately whether it is able to apply the frame, and applying it twice is harmless. The receiver
acknowledges \textbf{every fragment}, reporting the frame and fragment, its current decoder state hash, and its current estimate of the packet
inter-arrival time.

\subsection{Measuring capacity at the receiver, with a grace period}
\paperfig[0.55]{sf_fig2_grace}{The receiver averages the gaps $t_i$ between arriving packets. Because the sender pauses between frames, each fragment carries
a ``grace period'' (how long the sender itself waited before sending it), which the receiver subtracts from the measured gap.}

As in Sprout, the receiver estimates the bottleneck's speed from how closely packets arrive back to back: if fragments are queued at a bottleneck
that sends one packet per millisecond, they come out 1\,ms apart. It keeps a running average $\tau$ of the gap.

\textbf{The catch for video.} Bulk transfers send continuously, but a video sender sends one frame's fragments and then goes quiet until the next frame.
The gap between the last fragment of one frame and the first of the next is mostly the \emph{sender's} pause, not the network's speed. Counting it
would make the network look much slower than it is.

\textbf{The fix.} Each fragment carries a \textbf{grace period}: how long the sender waited between sending this fragment and the previous one. The receiver
removes it:
\[
  \tau_i \leftarrow \alpha\,(T_i - T_{i-1} - \text{grace}_i) + (1 - \alpha)\,\tau_{i-1},\qquad \alpha = 0.1,
\]
where $T_i$ is the arrival time of fragment $i$. With $\alpha = 0.1$ this behaves roughly like an average over the last 10 packets.

\textbf{Example.} Fragments arrive 1\,ms apart within a frame, so $\tau = 1$\,ms. Then the sender pauses 29\,ms before the next frame, which arrives 30\,ms
after the previous fragment.
\begin{center}
\begin{tabular}{@{}lll@{}}
\toprule
 & Gap used & New $\tau$\\
\midrule
without grace period & 30\,ms & $0.1 \times 30 + 0.9 \times 1 = 3.9$\,ms: the link looks 4$\times$ slower\\
with grace period & $30 - 29 = 1$\,ms & $0.1 \times 1 + 0.9 \times 1 = 1$\,ms: correct\\
\bottomrule
\end{tabular}
\end{center}
When the authors removed the grace period, Salsify underestimated capacity and sent low-quality video.

\subsection{How big may the next frame be?}
The sender wants every frame to arrive within $d = 100$\,ms, short enough to feel interactive. If the bottleneck delivers one packet every $\tau$, then
$d / \tau$ packets can be in the pipe without anything waiting longer than $d$. Some of that room is already used by packets sent but not yet
acknowledged:
\[
  N = (\text{index of last packet sent}) - (\text{index of last packet acknowledged}),
\]
so the next frame may use what remains:
\[
  \texttt{max\_frame\_size} = \Big(\frac{d}{\tau} - N\Big) \times \text{MTU}.
\]

\textbf{Example.} $\tau = 1$\,ms (about 12\,Mb/s with 1500-byte packets), $N = 40$ packets in flight. Allowed in flight: $100 / 1 = 100$ packets. Room
for this frame: $100 - 40 = 60$ packets $= 90$\,KB. The higher-quality version came out at 110\,KB (too big) and the lower at 70\,KB (fits), so the
lower one is sent, and the next frame's two candidates bracket that lower quality. If both versions had been over 90\,KB, the frame would be skipped and
the next camera picture encoded from the same state.

\subsection{The complete sender}
\begin{lstlisting}
on receiving an ACK:
  update mean_interarrival_time, known_receiver_codec_state,
         num_packets_outstanding
  if the ACK shows a loss: enter loss-recovery mode for 5 seconds
  max_frame_size = MTU * (100 ms / mean_interarrival_time
                          - num_packets_outstanding)

for each new camera image:
  source = loss_recovery_mode ? known_receiver_codec_state
                              : last_sent_frame.target_state
  lower  = encode(source, image, last_sent_frame.quality - DECR)
  higher = encode(source, image, last_sent_frame.quality + INCR)
  if higher.length < max_frame_size:            send higher
  elif lower.length < max_frame_size
       or consecutive_skips >= 4:               send lower
  else:                                         skip; consecutive_skips += 1
  (whenever a frame is sent: consecutive_skips = 0, last_sent_frame = it)
\end{lstlisting}
\textbf{Why send the lower frame after 4 skips even if it does not fit?} If Salsify skipped forever, nothing would be in flight, no ACKs would come back,
and the capacity estimate would never update. Salsify would never discover that the network had recovered. Sending something keeps the feedback loop
alive.

% ------------------------------------------------------------------
\section{Why both halves are needed}
The authors removed one piece at a time:
\begin{itemize}
  \item \textbf{Without the grace period}, capacity is underestimated, and quality and bitrate drop.
  \item \textbf{With a conventional codec} (state hidden, so only one encoding per frame, its quality chosen beforehand by trial-compressing a downscaled
        copy to aim at the target size), results were much worse. The transport has no choice at sending time: if the capacity estimate changes while the
        frame is being compressed, Salsify either adds delay or misses a chance to send better quality. And predicting a frame's size in advance is
        inherently inaccurate.
  \item With both removed, the system behaves like conventional videoconferencing.
\end{itemize}

\section{How it was measured (in brief)}
The testbed treats each video system as a black box. It appears to the sender as a USB webcam playing a repeatable 720p clip, and it captures the receiver's
HDMI output. Between them, a network emulator replays recorded cellular traces. Every frame carries barcodes so that each displayed frame can be matched to the
camera frame it came from. \textbf{Delay} is display time minus capture time; a frame that is never shown is charged until the next frame appears, so dropped
frames count against a system. \textbf{Quality} is SSIM against the original. On an emulated AT\&T LTE trace, Salsify-1c had lower 95th-percentile delay than
FaceTime, Hangouts, Skype and WebRTC (WebRTC's was 10.5$\times$ higher) and higher quality than all of them.

\section{Limitations}
\begin{itemize}
  \item \textbf{No audio.} Adding it raises open questions (audio usually sits in a jitter buffer; should video wait to stay in sync?).
  \item \textbf{Most codecs cannot save and restore state}, including the hardware encoders that phones need to save power.
  \item A conventional codec that could reliably hit a target frame size would reduce Salsify's advantage.
  \item The gains are largest on \textbf{highly variable} paths, such as cellular on the move. On steady networks Salsify is close to existing systems.
  \item No end-to-end scheme can keep delay low if the bottleneck queue is shared with other traffic that keeps it full.
\end{itemize}

\section{Salsify vs.\ Pensieve: two kinds of video}
\begin{center}
\begin{tabularx}{\linewidth}{@{}lXX@{}}
\toprule
 & \textbf{Pensieve (on-demand streaming)} & \textbf{Salsify (real-time video)}\\
\midrule
Video & pre-encoded chunks at fixed bitrates & encoded live, frame by frame\\
Decision & bitrate of the next chunk (seconds of video) & size and quality of the next frame (tens of ms)\\
Delay budget & seconds of buffer & about 100\,ms end to end\\
Transport & HTTP/TCP from a CDN & own UDP protocol, merged with the codec\\
How it adapts & learned policy over observations & measures capacity per frame, encodes to fit\\
Main danger & rebuffering & loss and queueing caused by its own frames\\
\bottomrule
\end{tabularx}
\end{center}

\begin{keybox}[Exam one-liner]
Salsify merges codec and transport rate control into one per-frame loop. The receiver estimates packet inter-arrival time with the sender's grace period
removed, and the sender sets \texttt{max\_frame\_size} $= (100\,\text{ms}/\tau - N) \times$ MTU. A purely functional VP8 codec encodes each frame at
one quality step above and one below the last, from the same saved state; Salsify sends the larger one that fits, else skips (sending the lower one after
4 skips). After a loss, it encodes from the state the receiver has acknowledged, instead of retransmitting.
\end{keybox}

\section{Revision checklist}
\begin{itemize}
  \item Codec state = reference images + probability tables; encoder and decoder must stay in step; a normal encoder cannot undo an encode or hit an exact size.
  \item Problem: two loops, average rate, 1\,s updates; oversized frames sent anyway when capacity drops (10 $\to$ 2\,Mb/s example).
  \item Target: access/cellular links where your own traffic causes loss and delay; FEC weak against bursts.
  \item Loop: budget $\to$ encode at last quality $\pm$ one step from same state $\to$ send higher if it fits, else lower, else skip; 4 skips $\Rightarrow$ send lower.
  \item Functional codec: \texttt{decode(state, frame)}, \texttt{encode(state, image, quality)}; try several, skip after encoding, choose the source state.
  \item Loss recovery: ACK carries decoder state hash; encode from the known receiver state for 5\,s (F3 lost $\Rightarrow$ F6 built on $S_2$); no retransmission.
  \item Transport: UDP fragments with source/target state hashes; ACK every fragment.
  \item Estimate: $\tau \leftarrow 0.1(T_i - T_{i-1} - \text{grace}_i) + 0.9\tau$ (30\,ms gap with 29\,ms pause counts as 1\,ms).
  \item Budget: $(d/\tau - N) \times$ MTU, $d = 100$\,ms (example: 60 packets = 90\,KB).
  \item Ablation: no grace period underestimates; a conventional codec gives no choice at send time.
  \item Limits: no audio, codecs (especially hardware) lack save/restore, gains mainly on variable paths.
\end{itemize}
